% ===========================================================================
% chapters/conclusions.tex — Conclusions
% ===========================================================================

\chapter*{Conclusions}
\addcontentsline{toc}{chapter}{Conclusions}
\markboth{CONCLUSIONS}{CONCLUSIONS}

This thesis has developed and exercised a dynamic mathematical model of a non-catalytic tubular reformer for pyrolytic gas, and has benchmarked three time-integration algorithms on the resulting stiff convection--diffusion--reaction system. The work delivers three artefacts: a ten-species, eight-reaction lumped kinetic scheme calibrated for iso-thermal operation; a one-dimensional plug-flow reactor with staged side-stream injection, discretised by first-order upwind finite volumes; and an integrator recommendation that meets the spatial-discretisation accuracy floor at materially lower cost than the alternatives.

The engineering results show a reactor with a narrow but real design margin. The temperature lever exposes a tar-clearance threshold just below the design point: at \(T = \SI{1073}{\kelvin}\) benzene steam reforming is too slow to consume the aromatic load within the available residence time, and a non-trivial tar-slip arrives at the outlet. The feed-velocity lever reveals two competing failure modes at the extremes --- over-oxidation at low feed velocity and under-conversion at high --- with the design value sitting in the narrow window between them. The side-injection flow-rate variation is the only one of the three with a non-monotonic optimum inside the tested range, and the design point sits close to that optimum. It is therefore the most directly actionable knob for any future parametric optimisation.

The numerical results identify forward Euler at \(\Delta t = \SI{1}{\milli\second}\) as the cheapest integrator that meets the accuracy floor at the engineering-relevant grid. Backward Euler at a tuned \(\Delta t = \SI{100}{\milli\second}\) is faster still, at the cost of user-side stability tuning; DLSODE provides default-configuration robustness but is penalised on the present iso-thermal model by a BDF-order-collapse pathology at fine grids. The recommendation is therefore tied to the present kinetic scheme and operating point: a richer reaction mechanism or a coupled energy balance would tighten the explicit method's stability bound and shift the ranking toward the implicit alternatives. Both extensions are natural directions for follow-on work, and the parametric studies in this thesis define the operating range over which the resulting model should be tested.
